\section{Related Work}

\paragraph{Recursive quantitative research.}
AQuA retains validated experiments to improve later symbolic-factor and
trainable-model proposals under fixed research scaffolds
\citep{guo2026aqua}. Our target is complementary: the persistent evolutionary
object is the Worker harness, while quantitative invariants provide an
observation of predicted Research State changes in that harness.

\paragraph{Harness optimization.}
AHE evolves a decoupled coding-agent harness from layered trajectory evidence,
edit predictions, and next-round task deltas \citep{lin2026ahe}.
Meta-Harness gives an agentic proposer access to harness source, scores, and the
full execution history of prior candidates \citep{lee2026metaharness}.
Our claim is not that the outer harness-evolution loop is new. We study whether
an explicit expected--observed--target Quant Research State transition, tested
through component activation and quantitative observations, improves search
and experience reuse. AHE is evaluated through a separately pinned,
source-faithful reproduction with a minimal quantitative-task adapter; it is
not reimplemented as a treatment switch inside the QRS runner. A QRS-no-State
ablation addresses quant-specific mechanism attribution, while the independent
AHE reproduction provides the external harness-optimization comparison.

\paragraph{Domain-grounded verification and adaptive harnesses.}
GLEAN compiles domain guidelines into trajectory-informed calibrated signals and
uses uncertainty to trigger active differential checks
\citep{zhang2026glean}. CHILL-Harness learns counterfactual orchestration
adaptations for long-horizon agents \citep{fu2026chill}. We do not claim to
invent domain-grounded verification or causal harness intervention. Our focus is
an open quantitative Research State representation for selecting persistent
full-harness components; answer-free quantitative invariants are used as
transition observations when applicable.
